\begin{abstract}
\noindent
Two items occur in both the \emph{Kalkulus I} and the \emph{Line\'aris algebra I}
\emph{t\'argytematika}: the \textbf{triangle inequality} (``nevezetes
egyenl\H otlens\'egek'' / ``vektorok hossza, h\'aromsz\"og-egyenl\H otlens\'eg'') and the
\textbf{inverse with respect to composition} (``inverz f\"uggv\'eny, \"osszet\'etel'' /
``m\'atrixok inverze''). The companion Lean files show that in each case the two courses
state one thing: the triangle inequality \emph{is} composition in a category enriched
over $([0,\infty),\ge,+)$, and the inverse \emph{is} isomorphism in a category. This
document asks the complementary question: \emph{how do you draw that, and how far can a
drawing alone carry a proof?} We survey eight diagrammatic languages -- commutative
diagrams, string diagrams (``graphical linear algebra''), proofs without words,
number-line and $\varepsilon$-interval pictures, Cartesian graphs and the mirror $y=x$,
Euler/Venn and bag-and-arrow diagrams, order (Hasse) diagrams, and geometric optimisation
pictures. For each we state its marks, its legal moves and whether it is certified, and
then work through the exercises of \lean{kalkulus\_gyakorlo.pdf} that can be settled by
drawing. The last section is a playbook for a timed exam. Every mathematical claim
illustrated here is stated and proved in \lean{RequestProject/DiagramProofs.lean}; the
captions name the corresponding theorem.
\end{abstract}

\tableofcontents

% ==========================================================================
\section{What has to be true for a drawing to be a proof}
\label{sec:what}

A picture is a proof when it is a term of a formal language whose syntax happens to be
two-dimensional. Three separate things must be in place.

\begin{description}[leftmargin=2.2em,style=nextline]
\item[Syntax.] A fixed finite stock of marks (arrows, boxes, wires, dots, segments,
  shaded regions) and a rule saying which arrangements of marks are well formed. A doodle
  that is not well formed is not a statement at all, so it can be neither a proof nor a
  mistake -- it is noise.
\item[Legal moves.] A list of allowed transformations of one well-formed diagram into
  another: pasting two commuting cells, sliding a box along a wire, cutting a region and
  reassembling it, reflecting a graph in $y=x$.
\item[Soundness.] A theorem -- proved once, symbolically, \emph{outside} the picture
  language -- saying: if a diagram $D$ can be turned into $D'$ by legal moves, then the
  statement denoted by $D$ implies the statement denoted by $D'$.
\end{description}

The third point splits the languages below into two very different classes.

\begin{itemize}
\item \textbf{Certified languages.} Commutative diagrams and string diagrams come with a
  proved soundness, and often completeness, theorem. For string diagrams in a monoidal
  category, two diagrams are related by planar deformation exactly when the corresponding
  algebraic terms are equal (Joyal--Street). Here the drawing \emph{is} the term: you are
  not illustrating a proof, you are writing one.
\item \textbf{Heuristic languages.} A dissection picture, a number-line sketch or a
  Cartesian graph has no such theorem behind it. It proves something only relative to a
  small \emph{symbolic residue}: which configuration is generic, which case is missing,
  whether the curve really is monotone. In practice the residue is one or two lines.
\end{itemize}

\begin{center}
\small
\begin{tabular}{@{}llll@{}}
\toprule
Language & Marks & Legal move & Certified? \\
\midrule
Commutative diagrams   & objects, arrows        & pasting, filling & yes \\
String diagrams        & boxes, wires, spiders  & planar isotopy, rewriting & yes \\
Proofs without words   & regions, lengths       & dissection, superposition & no \\
Number line            & points, intervals      & translation, reflection & no \\
Cartesian graphs       & curves, lines          & reflection, transformation & no \\
Euler / bag-and-arrow  & blobs, dots, arrows    & containment, chaining & partly \\
Hasse / order          & nodes, upward edges    & path composition & yes \\
Geometric optimisation & figures, envelopes     & superposition & no \\
\bottomrule
\end{tabular}
\end{center}

\medskip
\noindent
Hence the rule of thumb we return to in Section~\ref{sec:exam}: \emph{use a certified
language when the statement is about composition or invertibility; use a heuristic
language when the statement is about a quantity, and write its symbolic residue on one
line next to the picture.}
